%!TEX program = xelatex
% 中文翻译，原文：Log GLSM and related topics (Patrick Lei)
% 本译文由 Muse Spark 1.3 生成，未参考其他中文译本，直接译自 notes.tex 英文原文。
\documentclass{amsart}
\usepackage{amsmath}
\usepackage{amssymb}
\usepackage{amsthm}
\newtheorem{thm}{定理}[section]
\newtheorem{cor}[thm]{推论}
\newtheorem{prop}[thm]{命题}
\newtheorem{lem}[thm]{引理}
\newtheorem{conj}[thm]{猜想}
\newtheorem{quest}[thm]{问题}
\newtheorem{claim}[thm]{断言}

\theoremstyle{definition}
\newtheorem{defn}[thm]{定义}
\newtheorem{defns}[thm]{定义}
\newtheorem{con}[thm]{构造}
\newtheorem{exm}[thm]{例}
\newtheorem{exms}[thm]{例}
\newtheorem{notn}[thm]{记号}
\newtheorem{notns}[thm]{记号}
\newtheorem{addm}[thm]{补遗}
\newtheorem{exer}[thm]{习题}

\theoremstyle{remark}
\newtheorem{rmk}[thm]{注记}
\newtheorem{rmks}[thm]{注记}
\newtheorem{warn}[thm]{警告}
\newtheorem{sch}[thm]{附注}


% unnumbered theorems
\theoremstyle{plain}
\newtheorem*{thm*}{定理}
\newtheorem*{prop*}{命题}
\newtheorem*{lem*}{引理}
\newtheorem*{cor*}{推论}
\newtheorem*{conj*}{猜想}

% unnumbered definitions
\theoremstyle{definition}
\newtheorem*{defn*}{定义}
\newtheorem*{exer*}{习题}
\newtheorem*{defns*}{定义}
\newtheorem*{con*}{构造}
\newtheorem*{exm*}{例}
\newtheorem*{exms*}{例}
\newtheorem*{notn*}{记号}
\newtheorem*{notns*}{记号}
\newtheorem*{addm*}{补遗}


\theoremstyle{remark}
\newtheorem*{rmk*}{注记}
%\usepackage{MnSymbol}
\usepackage{bm}
\usepackage{accents}
\usepackage{mathtools}
\usepackage{tikz}
\usetikzlibrary{calc}
\usetikzlibrary{decorations.pathmorphing,shapes}
\usetikzlibrary{automata,positioning}
\usepackage{tikz-cd}
\usepackage{forest}
\usepackage{braket} 
\usepackage{listings}
\usepackage{mdframed}
\usepackage{verbatim}
\usepackage{physics2}
\usephysicsmodule{ab,ab.legacy,diagmat,xmat}
\usepackage{derivative}
\usepackage{fixdif}
\usepackage{stmaryrd}
% \usepackage{euscript} 
% \usepackage[mathcal]{eucal}
\usepackage{stackengine} 
%\usepackage{/home/patrickl/homework/macaulay2}

%font
\usepackage[sc]{mathpazo}
\usepackage{inconsolata}
\usepackage{microtype}
% \usepackage{fontspec} 
% \setmainfont{Tex Gyre Pagella}
\usepackage[OT1,euler-digits]{eulervm}
% \usepackage{euler-math} 
\usepackage[scaled=0.86]{berasans}

%CS packages
\usepackage{algorithmicx}
\usepackage{algpseudocode}
\usepackage{algorithm}

% typeset and bib
\usepackage[english]{babel} 
\usepackage[UTF8, scheme=plain]{ctex}
\usepackage[backend=biber,style=alphabetic,maxalphanames=4,maxnames=5,hyperref,isbn=false]{biblatex}
\usepackage[bookmarks, colorlinks, breaklinks]{hyperref} 
\hypersetup{linkcolor=blue,citecolor=magenta,filecolor=black,urlcolor=blue}
\usepackage{cleveref}
\usepackage{graphicx}
\graphicspath{{./}}


% other formatting packages
\usepackage{float}
\usepackage{booktabs}
\usepackage[shortlabels]{enumitem}
\setitemize{noitemsep}
\usepackage{csquotes}
%\usepackage{titlesec}
%\usepackage{titling}
%\usepackage{fancyhdr}
%\usepackage{lastpage}
% \usepackage{parskip}
\newlist{mydescription}{description}{1}
\setlist[mydescription]{style=nextline,
                        font=\bfseries,
                        % Tweak the next 4 options as needed:
                        labelindent=1cm, 
                        leftmargin =2cm,
                        rightmargin=1cm,
                        topsep     =1ex
                       }
\setcounter{tocdepth}{1}
\usepackage{lipsum}

% delimiters
\DeclarePairedDelimiter{\gen}{\langle}{\rangle}
\DeclarePairedDelimiter{\floor}{\lfloor}{\rfloor}
\DeclarePairedDelimiter{\ceil}{\lceil}{\rceil}



% The theorem-like environments above all share the `thm' counter, so cleveref
% names every one of them ``Theorem'' and cannot format lists of them. Alias the
% counter back to each environment and give cleveref the names.
\usepackage{etoolbox}
\crefname{thm}{定理}{定理}
\Crefname{thm}{定理}{定理}
\AtBeginEnvironment{conj}{\crefalias{thm}{conj}}
\AtBeginEnvironment{prop}{\crefalias{thm}{prop}}
\AtBeginEnvironment{lem}{\crefalias{thm}{lem}}
\AtBeginEnvironment{cor}{\crefalias{thm}{cor}}
\crefname{conj}{猜想}{猜想}
\Crefname{conj}{猜想}{猜想}
\crefname{prop}{命题}{命题}
\Crefname{prop}{命题}{命题}
\crefname{lem}{引理}{引理}
\Crefname{lem}{引理}{引理}
\crefname{cor}{推论}{推论}
\Crefname{cor}{推论}{推论}

% shortcuts
\DeclareMathOperator{\Ima}{Im}
\newcommand{\A}{\mathbb{A}}
\newcommand{\G}{\mathbb{G}}
\newcommand{\N}{\mathbb{N}}
\newcommand{\R}{\mathbb{R}}
\newcommand{\C}{\mathbb{C}}
\newcommand{\Z}{\mathbb{Z}}
\newcommand{\Q}{\mathbb{Q}}
\newcommand{\E}{\mathbb{E}}
\newcommand{\V}{\mathbb{V}}
\renewcommand{\H}{\mathbb{H}}
\renewcommand{\k}{\Bbbk}
\renewcommand{\L}{\mathbb{L}}
\renewcommand{\P}{\mathbb{P}}
\newcommand{\M}{\mathcal{M}}
\newcommand{\Mbar}{\overline{\mathcal{M}}}
\newcommand{\g}{\mathfrak{g}}
\newcommand{\h}{\mathfrak{h}}
\newcommand{\n}{\mathfrak{n}}
\renewcommand{\b}{\mathfrak{b}}
\newcommand{\ep}{\varepsilon}
\newcommand*{\dt}[1]{%
   \accentset{\mbox{\Huge\bfseries .}}{#1}}
%\renewcommand{\abstractname}{Official Description}
\newcommand{\mc}[1]{\mathcal{#1}}
% \newcommand{\msc}[1]{\mathscr{#1}}
\newcommand{\T}{\mathbb{T}}
\newcommand{\mf}[1]{\mathfrak{#1}}
\newcommand{\mbf}[1]{\mathbf{#1}}
\newcommand{\bv}{\mbf{v}}
\newcommand{\bq}{\mbf{q}}
\newcommand{\bp}{\mbf{p}}
\newcommand{\btau}{\bm{\tau}}
\newcommand{\mr}[1]{\mathrm{#1}}
\newcommand{\on}[1]{\operatorname{#1}}
\newcommand{\ms}[1]{\mathsf{#1}}
\newcommand{\mt}[1]{\mathtt{#1}}
\newcommand{\ol}[1]{\overline{#1}}
\newcommand{\ul}[1]{\underline{#1}}
\newcommand{\wt}[1]{\widetilde{#1}}
\newcommand{\wh}[1]{\widehat{#1}}
\renewcommand{\div}{\operatorname{div}}
\newcommand{\1}{\mathbf{1}}
\newcommand{\2}{\mathbf{2}}
\newcommand{\3}{\mathbf{3}}
\newcommand{\I}{\mathrm{I}}
\newcommand{\II}{\mr{I}\hspace{-1.3pt}\mr{I}}
\newcommand{\III}{\mr{I}\hspace{-1.3pt}\mr{I}\hspace{-1.3pt}\mr{I}}
\renewcommand{\v}{\mbf{v}}
\newcommand{\w}{\mbf{w}}
\newcommand{\bmu}{\bm{\mu}}
\newcommand{\pre}{\mr{pre}}
\newcommand{\vir}{\mr{vir}}
\newcommand{\fl}{\mr{fl}}
\newcommand{\ps}[1]{\llbracket #1 \rrbracket}
\newcommand{\ls}[1]{\llparenthesis #1 \rrparenthesis}
\newcommand{\loc}{\mr{loc}}

\DeclareMathOperator{\Der}{Der}
\DeclareMathOperator{\Tor}{Tor}
\DeclareMathOperator{\Hom}{Hom}
\DeclareMathOperator{\End}{End}
\DeclareMathOperator{\Ext}{Ext}
\DeclareMathOperator{\ad}{ad}
\DeclareMathOperator{\Aut}{Aut}
\DeclareMathOperator{\Rad}{Rad}
\DeclareMathOperator{\Pic}{Pic}
\DeclareMathOperator{\NS}{NS}
\DeclareMathOperator{\supp}{supp}
\DeclareMathOperator{\Supp}{Supp}
\DeclareMathOperator{\depth}{depth}
\DeclareMathOperator{\sgn}{sgn}
\DeclareMathOperator{\spec}{Spec}
\DeclareMathOperator{\Spec}{Spec}
\DeclareMathOperator{\proj}{Proj}
\DeclareMathOperator{\Proj}{Proj}
\DeclareMathOperator{\ord}{ord}
\DeclareMathOperator{\Div}{Div}
\DeclareMathOperator{\Bl}{Bl}
\DeclareMathOperator{\coker}{coker}
\DeclareMathOperator{\ev}{ev}
\DeclareMathOperator{\st}{st}
\DeclareMathOperator{\pr}{pr}
\DeclareMathOperator{\ch}{ch}
\DeclareMathOperator{\Cont}{Cont}

\addbibresource{references.bib}

